\centerline{\bf \Huge Муниципальный листочек}
\centerline{\bf \Huge Геометрия + теория чисел}

\problem{1}
Дан прямоугольник $A B C D$. Прямая, проходящая через вершину $A$ и точку $K$ на стороне $B C$, делит весь прямоугольник на две части, площадь одной из которых в 5 раз меньше площади другой. Найдите длину отрезка $K C$, если $A D=60$.

\problem{2}
В выпуклом четырехугольнике $A B C D$ сторона $C D$ видна с середины стороны под прямым углом. Докажите, что $A D+B C \geqslant C D$.

\problem{3}
Дан выпуклый четырёхугольник $A B C D, X$ - середина диагонали $A C$. Оказалось, что $C D \| B X$. Найдите $A D$, если известно, что $B X=3, B C=7, C D=6$.

\problem{4}
Имеется дробь $\frac{1}{3}$. Каждую секунду к её числителю прибавляется 1 , а к знаменателю прибавляется 7 . Восточное поверье гласит: в тот момент, когда получится дробь, сократимая на 11, наступит конец света. Докажите, что не следует бояться наступления конца света.

\problem{5}
Сумма квадратов $n$ простых чисел, каждое из которых больше 5 , делится на 6 . Докажите что и $n$ делится на 6 .

\problem{6}
На доске написано число 12. Каждую минуту число умножают или делят либо на 2 , либо на 3 и результат записывают на доске вместо исходного числа. Докажите, что число, которое будет написано на доске ровно через час, не будет равно 54.